\clearpage
\section{What optimizing the finite start can achieve}\label{start:section}
The preceding minimax criterion uses chains whose initial point tends to one.
One might instead optimize that initial point, including the corresponding
localization loss. The following identity rules out an improvement by this
particular modification; it does not rule out stronger Fourier estimates.

For a 1-Lipschitz function \(g:[0,1]\to\mathbb R\), retain the notation
\(\Phi_\zeta(g)\) for the infimum of admissible chain costs from
\(1-\zeta\) to zero, and \(\Phi_0(g)=\lim_{\zeta\downarrow0}\Phi_\zeta(g)\).
Define
\[
 I_\zeta(g)=\zeta+g(1-\zeta)-g(1),\qquad 0<\zeta<1.
\]
The Lipschitz bound gives \(0\le I_\zeta(g)\le2\zeta\).

\begin{proposition}[Finite-start identity]\label{start:identity}
For every such profile,
\[
 0\le\Phi_0(g)-\Phi_\zeta(g)
 \le\operatorname{Var}^-(g;[1-\zeta,1])\le\tfrac12 I_\zeta(g).
\]
Consequently
\[
 \inf_{0<\zeta<1}\bigl(\Phi_\zeta(g)+I_\zeta(g)\bigr)=\Phi_0(g).
\]
\end{proposition}
\begin{proof}
For \(0<n<n'<1\), truncating a chain from \(n'\) when it crosses
\(n\) preserves admissibility and decreases the crossing edge's cost.
Thus the optimal cost is nondecreasing with its initial point. Conversely,
successively replacing a current point \(x>n\) by \(\max\{n,2x-1\}\)
gives a finite admissible bridge from \(n'\) to \(n\): before the last
step, \(1-x\) doubles. Each edge costs at most the negative variation
on its interval. Appending an arbitrarily near-minimizing chain from
\(n\) bounds the difference of optimal costs by
\(\operatorname{Var}^-(g;[n,n'])\). This proves monotonicity,
boundedness, and the first comparison on letting \(n'\uparrow1\).

A Lipschitz function is absolutely continuous and \(-1\le g'\le1\)
almost everywhere. Since \((-v)_+\le(1-v)/2\) on that interval,
\[
 \operatorname{Var}^-(g;[n,1])
 =\int_n^1(-g')_+
 \le\frac12\int_n^1(1-g')
 =\frac{1-n+g(n)-g(1)}2.
\]
It follows that \(\Phi_\zeta+I_\zeta\ge\Phi_0+I_\zeta/2\ge\Phi_0\).
Letting \(\zeta\downarrow0\) proves the reverse inequality for the infimum.
\end{proof}

Here \(I_\zeta\) is precisely the leading profile correction obtainable
from the initial local-constancy argument. To see this, write \(f(x)=x+g(x)\).
At terminal depth \(N\), suppose occupied terminal cubes have mass at most
\(2^{-f_N(N)}\), and an initial cube of depth \(n\) has mass at least
\(2^{-f_N(n)-eN}\). Localization at bandwidth \(2^N\) contributes
\(2^{2N-f_N(N)}\). Converting the Lebesgue integral on the initial cube
to its normalized average contributes \(2^{-2n}\), and division by its
mass contributes at most \(2^{f_N(n)+eN}\). The resulting exponent is
\[
 2(N-n)+f_N(n)-f_N(N)+eN.
\]
For \(n=(1-\zeta)N\), its leading normalized part is
\(2\zeta+f(1-\zeta)-f(1)=I_\zeta(g)\). Enlargements, grid errors,
and regularization losses must still be added. Proposition~\ref{start:identity}
therefore applies even when this exact mass-profile correction replaces
the coarser loss \(2\zeta\). It is a limitation of that estimate, not
a lower bound for all possible localization operators.

\subsection{Why a zero shell exponent alone is insufficient}
There are compact null-length subsets of the line with positive probability
approximations whose mutually orthogonal quadratic increments are uniformly
bounded. This elementary example clarifies an endpoint issue; it is not a
planar distance-set counterexample.

Start with \(E_0=[0,1]\). At step \(j\ge1\), subdivide every surviving
interval into \(j+1\) equal pieces and retain the first \(j\). Use closed
intervals for the compact sets, ignoring endpoints in integrals. There are
\(j!\) intervals at stage \(j\), each of length
\(\ell_j=1/(j+1)!\), and \(|E_j|=1/(j+1)\). Put
\(p_j=(j+1)\mathbf1_{E_j}\). These are probability densities and preserve
the mass of each earlier interval. Their weak limit is a probability
\(\eta\) on \(E_\infty=\bigcap_j E_j\), which has length zero.

Direct computation gives
\[
 \|p_j\|_2^2=j+1,\qquad \int p_jp_{j-1}=j,
 \qquad \|p_j-p_{j-1}\|_2^2=1.
\]
The increments are mutually orthogonal: each has integral zero on every
earlier surviving interval, and the preceding increments are constant
there. The limiting support even has Hausdorff dimension one. Indeed,
each stage-\(j\) interval has mass \(1/j!\). If
\(\ell_j\le r<\ell_{j-1}\), an interval \(I\) of length \(r\) meets at most
\(r/\ell_j+2\) stage-\(j\) grid cells, giving
\[
 \eta(I)\le3(j+1)r.
\]
For each fixed \(0<t<1\),
\((j+1)\ell_{j-1}^{1-t}=(j+1)(j!)^{-(1-t)}\) is bounded.
Thus \(\eta(I)\le C_t r^t\), proving the dimension assertion.

Uniform bounds on individual quadratic increments therefore do not
establish positive support length. These increments are not asserted to
be Fourier annuli of a distance measure. An endpoint distance argument
still needs stronger information, such as a uniformly bounded full
positive approximation or another suitable nonconcentration estimate.
